% Standalone LaTeX edition of README.md; keep the two guides in sync.
% Build from the repository root:
% latexmk -norc -r examples/latexmkrc -lualatex examples/README.tex
\documentclass[11pt]{article}
\usepackage[a4paper,margin=25mm]{geometry}
\usepackage{fontspec}
\setmainfont{Latin Modern Roman}
\setmonofont{Latin Modern Mono}[Scale=MatchLowercase]
\usepackage{microtype}
\usepackage{booktabs}
\usepackage{tabularx}
\usepackage{fvextra}
\usepackage{xurl}
\usepackage[colorlinks=true,linkcolor=blue,urlcolor=blue]{hyperref}
\DefineVerbatimEnvironment{verbatim}{Verbatim}{fontsize=\small,breaklines=true,breakanywhere=true}
\newcommand{\code}[1]{{\mdseries\ttfamily\detokenize{#1}}}
\hypersetup{pdftitle={tex-lint: LaTeX User Guide}}
\title{\code{tex-lint}: LaTeX User Guide}
\author{}
\date{}

\begin{document}
\maketitle

\begin{abstract}
This is the standalone LaTeX edition of the project's Markdown README,
and its source is itself a demonstration of semantic line wrapping.
The prose is formatted with \code{tex-lint} using the project's default soft and
hard widths of 80 and 100 source columns.
Code examples remain verbatim, while headings, lists, tables, hyperlinks,
and inline commands demonstrate the handling of ordinary LaTeX structure.
\end{abstract}

\tableofcontents

\section{Overview}

Semantic source wrapping for LaTeX uses \code{latexindent.pl} for sentence
detection and protected regions.
The CLI and wrapping policy follow the companion
\href{https://github.com/nickfraser/mdformat-semantic-wrap}{\code{mdformat-semantic-wrap}}
project.

Each detected sentence starts on its own source line.
Short sentences stay together;
longer sentences prefer punctuation near a soft width,
with a separate hard width.
Widths count source display columns, including LaTeX markup and indentation,
rather than typeset text.

\section{Requirements and current status}

\begin{itemize}
\item Perl 5.32 or newer,
including its standard modules and Unicode character data.
\item The \textbf{Perl-script distribution} of \code{latexindent},
with its \code{LatexIndent/} modules and \code{defaultSettings.yaml}.
TeX Live's script distribution is suitable.
\item The Perl dependencies of \code{latexindent},
including \code{YAML::Tiny} and \code{File::HomeDir},
available to the Perl interpreter running \code{tex-lint}.
\end{itemize}

There is no project package-install or build step,
and no Python environment is needed.
With the runtime requirements available,
make \code{tex-lint} available on \code{PATH} as described below.

\textbf{Verified integration:} \code{latexindent} \textbf{3.24.7},
\code{latexmk} \textbf{4.87}, TeX Live \textbf{2025} (pdfLaTeX and LuaLaTeX),
and Poppler \textbf{25.07.0}.
The full suite passes \textbf{376 assertions},
including 14 original/formatted document pairs with identical extracted text,
page geometry, and page rasters at 96 dpi.
Tests also verify source idempotence and resolved bibliography/cross-references.

The adapter uses process-local overrides of the internal sentence methods of
\code{latexindent}.
The tested 3.x path handles the older settings hash and document-level
sentence-wrapping callback.
A separate 4.x path follows the internal interface of upstream 4.0.3;
it has callback unit tests but has not been verified against a real 4.x
installation.
The adapter checks the methods it needs and does not modify the installed tool.
A standalone compiled \code{latexindent} binary without its Perl modules is not
a supported backend.

\section{Install tex-lint on PATH}

From the repository root, create a symlink in your user-local executable directory:

\begin{verbatim}
mkdir -p "$HOME/.local/bin"
ln -s "$PWD/bin/tex-lint" "$HOME/.local/bin/tex-lint"
export PATH="$HOME/.local/bin:$PATH"
tex-lint --help
\end{verbatim}

Add the \code{export PATH=...} line to your shell's startup file to persist it
across sessions.
The symlink points to this checkout,
allowing the command to locate its \code{lib/} and \code{libexec/} supporting
files.
Keep the checkout at that location, or recreate the symlink after moving it.
The examples below assume \code{tex-lint} is on \code{PATH}.

\section{Locate TeX Live}

Once installation under \code{~/.local/texlive} is complete,
add the installation's actual platform-specific \code{bin} directory to
\code{PATH}.
Then check the backend:

\begin{verbatim}
latexindent --version
tex-lint --help
\end{verbatim}

Alternatively, select the script explicitly:

\begin{verbatim}
tex-lint --latexindent /absolute/path/to/latexindent.pl document.tex
\end{verbatim}

Setting \code{LATEXINDENT=/absolute/path/to/latexindent.pl} provides the same
selection without repeating the option.
The formatter resolves executable symlinks so the adjacent upstream Perl modules
and settings can be loaded.
The backend runs under the same Perl interpreter as the formatter;
a successful standalone version check under a different interpreter is not
sufficient.

\section{Format and check}

Like \code{mdformat-semantic-wrap}, file inputs are \textbf{formatted in place by default}:

\begin{verbatim}
tex-lint document.tex
tex-lint chapters/ appendix.tex
tex-lint --check document.tex
\end{verbatim}

Directories are searched recursively for \code{.tex} files in deterministic
order.
Hidden directories, \code{venv}, \code{env},
and \code{node_modules} are skipped;
directory symlinks encountered during recursion are not followed.
Explicit file inputs are accepted regardless of extension.
Duplicate resolved paths are processed once.
For an explicitly supplied directory symlink, its contents are searched.

Use a separate output file or stdout to preview:

\begin{verbatim}
tex-lint --output document-formatted.tex document.tex
tex-lint --output - document.tex
tex-lint - < document.tex > document-formatted.tex
\end{verbatim}

The \code{-} input must be the only input;
\code{--output} requires a single resolved input file.
Input and output text use UTF-8.
Formatting is completed for every input before any source file is edited.
Backend errors leave the inputs untouched.
Existing file permissions are preserved when writing in place.

\begin{center}
    \begin{tabularx}{\linewidth}{@{}lX@{}}
\toprule
Status & Meaning \\
\midrule
\code{0} & Formatting succeeded, or every input already passes \code{--check}. \\
\code{1} & \code{--check} found unformatted input. No files were changed. \\
\code{2} & Invalid arguments/configuration, unavailable backend, or another error. \\
\bottomrule
\end{tabularx}
\end{center}

Check-mode reports go to stderr.
The command \code{--check -} compares stdin without emitting formatted text.
The options \code{--check} and \code{--output} cannot be combined.

\section{LaTeX README demonstration}

The source of this guide is \code{examples/README.tex},
a maintained LaTeX counterpart of \code{README.md}.
It demonstrates the formatter on a substantial user guide, including prose,
lists, a table, hyperlinks, inline commands, and protected code examples.
Build it with LuaLaTeX for Unicode punctuation and verbatim configuration
examples.

Run these commands from the repository root:

\begin{verbatim}
tex-lint --check examples/README.tex
latexmk -norc -r examples/latexmkrc -lualatex examples/README.tex
\end{verbatim}

The PDF is written to \code{.ocs/examples/README.pdf}.
The example's \code{latexmkrc} also confines auxiliary files and TeX/font caches
to \code{.ocs/examples/}.
To preview the formatter's effect on a copy, use:

\begin{verbatim}
tex-lint --output .ocs/examples/README-formatted.tex examples/README.tex
\end{verbatim}

After a normal build, the output directory already exists.
The source is formatted at the project's default 80/100 widths;
the compilation suite additionally reformats a copy at 40/55 and verifies that
its PDF remains identical in text, geometry, and rendered pixels.

\section{Configuration}

Put \code{.latexindent.yaml} in the LaTeX project.
The formatter uses the nearest ancestor configuration for each input,
or the nearest configuration from the working directory for stdin.
The option \code{--config PATH} selects one file for every input.
Only that settings file and the defaults of \code{latexindent} are loaded;
ambient user configuration is excluded.
YAML \code{paths} includes are not supported.

\begin{verbatim}
semantic_wrap:
  soft_width: 80
  hard_width: 100
  preferred_markers: ',;:-–—'
  continuation_indent: 0

defaultIndent: '    '
\end{verbatim}

The checkout's configuration also contains the baseline \code{modifyLineBreaks}
settings from the original design note, \code{latexindent_sentence_wrapping.md}.
Running ordinary \code{latexindent -m -l} uses the baseline width wrapping;
the conditional punctuation policy requires \code{tex-lint}.

CLI options override the corresponding \code{semantic_wrap} settings:

\begin{verbatim}
tex-lint \
  --semantic-wrap-soft-width 72 \
  --semantic-wrap-hard-width 88 \
  --semantic-wrap-preferred-markers ',;:' \
  document.tex
\end{verbatim}

Widths must be positive integers and \code{soft_width} must not exceed
\code{hard_width}.
The setting \code{continuation_indent} is a nonnegative number of extra spaces,
also exposed as \code{--semantic-wrap-continuation-indent}.
It is additional to the surrounding LaTeX indentation.
With no project setting,
surrounding indentation follows the defaults of \code{latexindent},
including tab indentation.

\subsection{List source layout}

This project's configuration keeps the first sentence on the same line as \verb|\item|, and makes ordinary and nested lists flush-left without a hanging indent:

\begin{verbatim}
\begin{itemize}
\item First sentence.
Another sentence in the same item.
\item A longer sentence, with a preferred break,
and a continuation without extra indentation.
\end{itemize}
\end{verbatim}

Select that style through native latexindent settings:

\begin{verbatim}
noAdditionalIndent:
  item: 1
  itemize: 1
  enumerate: 1
  description: 1

noAdditionalIndentGlobal:
  items: 1

modifyLineBreaks:
  items:
    ItemFinishesWithLineBreak: -1
\end{verbatim}

The explicit \code{noAdditionalIndent.item: 1} is important:
upstream's per-name six-space item indent takes precedence over
\code{noAdditionalIndentGlobal.items}.
The environment entries remove the additional indentation inside those lists;
omit those entries to retain surrounding list indentation with inline items.
Any indentation inherited from a containing non-list environment still applies.
The setting \code{continuation_indent} remains available if extra continuation
spaces are desired.

The adapter honors \code{ItemFinishesWithLineBreak: -1} for standard \verb|\item| commands,
including balanced optional labels such as \verb|\item[Label]|.
The item command and label count toward the first line's width only;
the first word stays attached to the marker.
Later sentences still start on new source lines.
Comments, blank paragraphs,
and protected block content remain barriers rather than being joined onto an
item line.
A long indivisible marker/label plus its first word may overflow.

Without the \code{-1} setting, existing separate-item behavior is retained.

Set \code{ItemFinishesWithLineBreak: 1} to explicitly request a newline after
\verb|\item|.

\subsection{Wrapping policy}

While the remaining sentence exceeds the available soft width, choose:

\begin{enumerate}
\item The rightmost preferred punctuation boundary at or before soft.
\item Otherwise, the first preferred boundary after soft but at or before hard.
\item Otherwise, the rightmost safe whitespace boundary at or before soft.
\item If the first indivisible word or protected span exceeds soft,
the first safe whitespace after that span.
\end{enumerate}

For example, with soft \code{20} and hard \code{26}:

\begin{verbatim}
One, two three, four five six seven.
\end{verbatim}

becomes:

\begin{verbatim}
One, two three,
four five six seven.
\end{verbatim}

Markers have equal priority; their order has no effect.
A marker is a candidate only when followed by breakable whitespace.
Hyphens, en dashes, and em dashes also need preceding whitespace,
preserving hyphenated words.
When \code{-} is enabled,
TeX's standalone \code{--} and \code{---} spellings are also preferred
boundaries.
An empty marker string disables punctuation preference:

\begin{verbatim}
tex-lint --semantic-wrap-preferred-markers '' document.tex
\end{verbatim}

Only ordinary spaces and tabs are breakable; nonbreaking spaces are not.
Widths account for wide characters, combining accents,
and emoji grapheme clusters.
Tabs use eight-column stops.
Individual words, URLs,
and protected macros may overflow the hard width rather than being split.

\subsection{LaTeX-specific behavior}

\begin{itemize}
\item Sentence recognition, including its abbreviation heuristics,
comes from \code{latexindent}.
Customize its \code{oneSentencePerLine} detection settings for the project.
Markdown's \code{--slw-*} language options have no direct LaTeX equivalent.
Prose upstream does not recognize as a sentence,
such as text without a sentence-ending marker,
receives normal indentation rather than reflow.
\item Existing eligible prose line breaks are rejoined before wrapping.
\item Verbatim environments and commands use upstream's protection routines.
\item Inline math (\verb|$...$|, \verb|\(...\)|), display math,
common equation environments, common table environments,
and \code{tikzpicture} are protected from reflow.
\item Macro invocations with balanced optional/mandatory arguments stay atomic,
including URLs and citations.
Their original source width is measured.
Sentences inside a macro argument are not split independently.
\item Comments, paragraph boundaries, explicit \verb|\\| breaks,
and multiline protected objects are barriers.
Sentences crossing a barrier are conservatively left out of reflow.
Comment text is not wrapped or moved by the sentence pass.
\item Document structure, item indentation,
and other normal indentation remain the responsibility of \code{latexindent}.
Custom verbatim environments and \verb|%\begin{noindent}| blocks can be
specified through its usual YAML settings.
\end{itemize}

The adapter enforces sentence manipulation, wrapping after indentation,
and overflow preservation, and disables the upstream comment-wrapping pass.
The setting \code{sentenceIndent} is managed through \code{continuation_indent}
instead.
YAML switches that enable replacement passes, upstream check mode,
or stdout logging are rejected because they conflict with the formatter's CLI
protocol.

\section{Tests}

The dependency-free TAP harness runs with the existing Perl installation:

\begin{verbatim}
perl tests/run.pl
\end{verbatim}

Tests cover companion-compatible punctuation selection, Unicode widths,
configuration validation, protected spans, indentation callbacks, idempotence,
and CLI errors.
Without a real backend,
the full-document integration group is explicitly skipped.

Run the real-backend suite using \code{latexindent} from \code{PATH}:

\begin{verbatim}
perl tests/run.pl --integration
\end{verbatim}

To select a particular backend instead:

\begin{verbatim}
TEX_LINT_LATEXINDENT=/absolute/path/to/latexindent.pl perl tests/run.pl
\end{verbatim}

That group exercises actual sentence detection, YAML/CLI precedence,
file/check/stdin/output modes, recursive input discovery,
protected TeX fixtures, and whole-document idempotence.
A backend or dependency failure fails the tests.

\subsection{Compilation and PDF equivalence}

Run the complete suite, including real-backend integration and compilation:

\begin{verbatim}
perl tests/run.pl --compile
\end{verbatim}

This requires \code{latexindent}, \code{latexmk}, \code{pdflatex},
\code{lualatex}, \code{bibtex}, \code{pdftotext}, \code{pdfinfo},
and \code{pdftoppm} on \code{PATH}.
On Fedora, the PDF utilities come from \code{poppler-utils}.
Explicit integration/compilation requests fail if a required tool is unavailable;
the default unit-test command skips those groups.
The variable \code{TEX_LINT_LATEXINDENT} still overrides backend selection with
\code{--compile}.

Seven fixtures in \code{tests/fixtures/} cover ordinary prose,
abbreviation heuristics, macro arguments, nested lists, tables, math, comments,
explicit line breaks, verbatim/custom newline-sensitive environments, Unicode,
footnotes, URLs, BibTeX citations, multipass cross-references,
and inline/flush-left items.
Six fixtures run with both pdfLaTeX and LuaLaTeX;
the \code{fontspec} Unicode fixture and this LaTeX README example run with
LuaLaTeX.

For each document/engine combination, the suite:

\begin{enumerate}
\item Confirms that formatting actually changes source and is idempotent.
\item Builds original and formatted documents using \code{latexmk -norc},
allowing bibliography and reference passes to converge.
\item Checks the final logs for unresolved references and missing glyphs.
\item Compares exact \code{pdftotext -layout} output and page count/geometry.
\item Compares every rendered page byte-for-byte at 96 dpi using
\code{pdftoppm}.
\end{enumerate}

Separate negative controls deliberately change content and shift text placement
to verify that the comparisons catch both kinds of difference.
PDF files themselves are not compared byte-for-byte because metadata and
source-dependent objects can differ without changing typesetting.

Runtime scratch input copies, output copies,
and upstream logs are confined to \path{$PWD/.ocs/tex-lint/run-*} and removed
after each invocation.
Tests use disposable directories under \path{.ocs/tests/}.
Compilation runs retain source copies, PDFs, extracted text, page rasters,
build logs, and local TeX/font caches under \path{.ocs/integration/compile-*};
the test output prints the exact directory.
These artifacts are ignored by Git and can be inspected after a failure or
success.

\section{Pre-commit}

Use a local system-language hook to run the command installed on \code{PATH}:

\begin{verbatim}
repos:
  - repo: local
    hooks:
      - id: tex-lint
        name: Semantic LaTeX wrapping
        entry: tex-lint
        language: system
        files: '\.tex$'
        exclude: ^tests/fixtures/
\end{verbatim}

The hook formats matching files in place and uses the current
\code{PATH}/\code{LATEXINDENT}.
Ensure \code{tex-lint} is on the \code{PATH} inherited by pre-commit,
including when invoked from an editor.
This repository's own hook uses \code{entry: bin/tex-lint} to exercise the
current checkout directly.
It excludes deliberately unformatted regression fixtures so their source layout
remains part of the test input.

\section{Source layout}

\begin{description}
\item[\code{bin/tex-lint}] Input discovery, CLI,
and the check/write/output workflow.
\item[\code{lib/TexLint/Config.pm}] Shared setting validation.
\item[\code{lib/TexLint/SentenceWrap.pm}] Independent display-column wrapping
core.
\item[\code{lib/TexLint/Protection.pm}] Balanced macro-argument protection.
\item[\code{lib/TexLint/LatexIndentAdapter.pm}] Isolated upstream sentence
integration.
\item[\code{libexec/tex-lint-worker.pl}] One fresh upstream process per input
file.
\item[\code{tests/}] Dependency-free tests and gated real-backend fixtures.
\item[\code{examples/README.tex}] This maintained LaTeX user guide and
formatting demonstration.
\end{description}

\end{document}
